\section{Methods}
\label{sec:methods}

\subsection{Puzzles as finite-domain constraint problems}

Every puzzle in this study is a finite-domain constraint satisfaction problem with three parts: variables with small domains, \emph{rules} of the genre that every instance shares, and \emph{clues} that define one instance and can be removed. A rule or clue is either an all-different constraint or a table constraint, given by a predicate over the values of a small scope of variables. We compile the problem to CNF with the direct encoding: one Boolean $x_{v,a}$ per variable $v$ and value $a$, exactly-one clauses per variable, pairwise at-most-one clauses for all-different (plus at-least-one clauses when the scope is a permutation, such as a Sudoku row), and one clause per forbidden tuple of a table constraint. Uniqueness is checked by solving, adding a clause that blocks the solution and solving again; the second call must be unsatisfiable \citep{yato2003complexity,lynce2006sudoku}. We use MiniSat~2.2 through PySAT \citep{een2003minisat,ignatiev2018pysat} and repeat every certificate with an independent model in OR-tools CP-SAT \citep{ortools}. Both solvers agreed on every puzzle we certified.

\paragraph{Families.} We use four families. \emph{Zebra} puzzles have $N$ houses and $K$ categories with $N$ values each (default $N{=}5$, $K{=}4$); clues are generated from 11 relation types (same person, different people, house position, not at a position, directly left of, next to, not next to, somewhere left of, exactly one house between, at an end, either--or) and rendered into English by templates with several phrasings each. Each zebra clue carries a canonical expression in the position DSL described below, which is the ground truth for its meaning. \emph{Sudoku} is the classic $9{\times}9$ game. \emph{Killer Sudoku} is played on $6{\times}6$ grids with $2{\times}3$ boxes; clues are cages (sum, digits distinct) and, rarely, givens. \emph{Futoshiki} is played on $5{\times}5$ Latin squares with inequality clues between neighbours and optional givens.

\paragraph{Certified generation.} For each family we draw a random solution, build a pool of clues that are true in it, reach uniqueness, and minimise. For zebra puzzles we start from the empty clue set and repeatedly add a pool clue that is violated by the current counter-example solution, which guarantees termination. For the grid families we start from the full pool. Minimisation removes clues one at a time in random order (givens first for Killer and Futoshiki) and keeps a removal whenever the solution stays unique, so every clue of a released puzzle is necessary. All randomness is seeded; every stage writes JSONL checkpoints and is resumable.

\subsection{The difficulty ladder}
\label{sec:ladder}

We grade a puzzle by the weakest deduction engine that reaches the full solution without search:
\begin{description}[leftmargin=1.4em,itemsep=1pt,topsep=2pt]
\item[L1] unit propagation on the CNF. In Sudoku terms this covers naked and hidden singles; for clues it covers using a clue once the other side is known.
\item[L2] L1 plus single-constraint elimination: generalised arc consistency on each rule or clue taken on its own, together with the exactly-one constraints of its scope \citep{mackworth1977consistency,regin1994alldiff}. This covers eliminations licensed by one clue, naked and hidden subsets inside one row or category, and cage combinations.
\item[L3] L2 plus failed-literal probing: place a candidate value, run unit propagation, and eliminate the candidate if this yields a contradiction \citep{li1997lookahead,heule2009lookahead}.
\item[L4] L2 plus two-level lookahead: place a candidate, then run unit propagation and failed-literal probing inside that hypothesis.
\item[L5] none of the above suffices; the solver must search.
\end{description}
All rules are sound and monotone, so each engine has a unique fixpoint, and the level can be computed in one pass that enables stronger rules only when weaker ones are exhausted. We call a puzzle \emph{fair} if it is unique and graded L1--L3. Unit tests check the ladder on hand-built instances that sit exactly at each of the five levels, on the classic Einstein puzzle, and on a published Sudoku with known answer.

\paragraph{Solver-side metrics.} For each puzzle we record (i) the fraction of variables fixed after the L1 and L2 fixpoints (a ``propagation backbone''; the full backbone of a unique puzzle is trivially every variable), (ii) MiniSat decisions, conflicts and propagations for finding the solution and proving it unique, (iii) for every non-given fact, the size of a minimal set of clues that forces it, computed as a group minimal unsatisfiable subset over clue selectors by deletion \citep{liffiton2008mus,bogaerts2021stepwise}, and (iv) clue redundancy, the fraction of clues whose individual removal keeps the solution unique.

\subsection{LLM calls}
\label{sec:llmcalls}

No API key was available in the environment where this study was run, and open-weight models could not be downloaded. Every LLM call was therefore executed as a fresh Claude Code subagent with no shared context, selected by model alias (Haiku, Sonnet, Opus); a probe call confirmed the alias resolved to \texttt{claude-haiku-5-5}, and the other aliases are documented as \texttt{claude-sonnet-5-5} and \texttt{claude-opus-5-5}. A pipeline stage writes one task file per call; the worker is instructed to read that file, reason without tools, and write one JSON output file. The harness reports the total token count and the number of tool calls of each subagent, which we log. Workers are told not to run code; in \NViolations{} of \NCalls{} calls the log shows one extra tool call, which in the cases we inspected was a read-only status query that the worker itself reported, not code execution. Several items share a call (10 puzzles per encoding or generation call, 4--5 per solving call), so items within a call are not independent of each other, but different encoders or solvers of the same item never share a call. Model reasoning settings were the harness defaults and are not under our control. The same task files can be run against the Messages API with one function in the repository.

\subsection{E1: dual-encoding verification}

An encoder receives the puzzle text and a short description of a position DSL in which \texttt{pos('tea')} is the house number of a value and each clue is one Python Boolean expression over \texttt{pos}, integers, comparisons, \texttt{Abs}, \texttt{And}, \texttt{Or}, \texttt{Not} and \texttt{Implies}. The prompt gives three worked examples (same person, fixed position, next to) and leaves every other relation to the model. The same expression can be evaluated concretely (to build table constraints) or symbolically in z3 \citep{demoura2008z3}. Two encodings $A$ and $B$ are equivalent when $\mathit{base} \wedge (A \oplus B)$ is unsatisfiable, where $\mathit{base}$ states the domains and the per-category all-different rules. The same test against the generator's own expressions gives a ground-truth check for single encodings, at the level of the whole puzzle and of each clue.

We encoded 299 zebra puzzles in three text conditions: 100 with plain templates; 100 with ``hard'' templates that use indirect references (``the left-hand neighbour of the Dane keeps the cat''), counting from the right, negated orderings, conditionals, ``neither \ldots nor'' and ``between''; and 99 whose plain-template clues Sonnet rewrote into natural prose under an instruction to keep the meaning exactly (one of 100 paraphrases was misaligned and excluded). Each puzzle was encoded three times: twice by Sonnet in independent calls with different batch compositions (A and B), and once by Haiku (C). For every clue-level disagreement with the ground truth we assigned an error category automatically, by testing equivalence with a list of plausible misreadings, and then read every case by hand.

\subsection{E2: validity of LLM-written puzzles}

Sonnet wrote 40 zebra, 20 Sudoku, 20 Killer and 20 Futoshiki puzzles together with their intended solutions, in batches of ten, under an instruction that each must have exactly one solution and be solvable without guessing. Zebra puzzles were dual-encoded as in E1 and kept only when the two encodings were equivalent. A puzzle is \emph{valid} if it parses, has exactly one solution, that solution is the stated one, and it is fair. The solver loop then (i) returns the solver's feedback (a second solution, the clues the stated solution violates, or a minimal conflicting clue set) to a fresh Sonnet call for repair, and (ii) minimises the puzzle by clue removal. We also measured redundancy and searched for repeated Sudoku layouts.

\subsection{E3: ladder, solver metrics and LLM accuracy}

From the certified banks (1,500 minimal puzzles per family, plus 8,000 extra Sudoku drawn to look for rare high levels) we sampled up to eight puzzles per family and level. Minimal zebra, Killer and Futoshiki puzzles are almost never L1, so the L1 cells were filled by adding random true clues to minimal L2--L3 puzzles until they reached L1; these are marked as augmented. No minimal Sudoku above L3 occurred among 9,500 generated, and only four L5 Killer puzzles occurred among 1,500. The final set has 116 puzzles. Haiku, Sonnet and Opus each solved every puzzle (4--5 per call, same batches for every model). We score exact solution and cell accuracy (over non-given cells); an unanswered item scores zero. We relate level to solver metrics by Spearman correlation with bootstrap intervals \citep{efron1993bootstrap}, and to accuracy with a logistic regression of exact solution on level, family and model, with puzzle-level cluster bootstrap for the level coefficient, and a linear mixed model of cell accuracy with a random intercept per puzzle \citep{seabold2010statsmodels}.

\subsection{E4: evolving fair puzzles against a solver}

We ran a $(\mu{+}\lambda)$ evolution of zebra puzzles with $\mu=8$. Each generation Haiku (the training solver) solves every individual; fitness is one minus the mean cell accuracy over all of an individual's evaluations, so survivors are re-evaluated every generation and lucky scores average out; ties prefer fewer clues. Sonnet proposes two variants per parent as edit operations on the formal clue list (add, remove or replace a clue of a typed relation) after seeing the puzzle, its solution and the solver's mistakes. The pipeline applies valid edits, discards clues that are false in the solution, restores uniqueness with random true clues if needed, and rejects any child that is not fair (L$>$3). Held-out testing re-solves the initial and final populations with fresh Haiku calls and with Sonnet and Opus.

\subsection{E5: LLM-invented genres}

Sonnet was asked to invent six new genres and to write each as Python code against a small API: a grid size, a function that adds global rules (all-different or table constraints over cells) and a function that, given a solution grid, lists candidate clues true of it with a checking predicate and a text. The pipeline compiles the code in a restricted namespace, draws random solutions of the rules with a SAT solver, builds a minimal unique clue set from the candidate clues, and grades the result. We generated twelve instances per genre.
